\section{The nonzero coefficient and the original knot volume}
\label{cta:conclusion}

\subsection{The coefficient of the whole invariant}
Let
\[
 \mathcal Q=-\Psi''(\boldsymbol\tau_*)-\ii\diag(b_v).
\]
Its real quadratic part is positive on real vectors.
Choose \(\sqrt{\det\mathcal Q}\) by the complex Gaussian on that
real tangent. This branch continues from positive definite real
matrices in the domain with positive real part.
Applying \eqref{cta:packet-expansion} to the first integral
in \eqref{cta:finite-charges} gives the explicit constant
\begin{equation}\label{cta:amplitude}
 D_{\mathcal T,\mathbf n}
 =\left(-\frac1{\pi\ii}\right)^r(2\pi)^{r/2}
   \frac{a(\boldsymbol\tau_*)\prod_v\sin\tau_{*,v}}
        {\sqrt{\det\mathcal Q}}\ne0 .
\end{equation}
Each factor has already been proved nonzero.
The charge-one coefficient is exactly one at every vertex,
even at a colliding or constant grid. Hence no parity-dependent
leading cancellation remains.

The powers of \(N\) can be read without an asymptotic
normalization choice. The \(r\) exact trace factors supply
\(N^{2r}\). The \(e\) finite edge sums supply \(N^{-3e/2}\),
and the angular Gaussian supplies \(N^{-r/2}\).
Since \(e=r-1\), the total is \(N^{3/2}\).
At \(s=0\) all physical turn factors and \(H_N/W_k\) are one.
Equations \eqref{cta:exact-source} and \eqref{cta:trace} therefore
leave precisely
\[
 \vartheta_d^{-w(D)}\prod_v C_{N,n_v}(0)
       =q_N^{d^2(2w(D)+\sum_vn_v)/4}=\Omega_N.
\]
Every final turn is included by the original compiler in
\eqref{cta:exact-source}; it is not removed in this count.
This proves the complete leading term \eqref{cta:leading}.
All nonminimal charges and fixed-loss errors are absorbed
using \eqref{cta:charge-order} and \eqref{cta:budget}.
The implied exponential rates and initial levels depend on
the fixed tree and powers, as the theorem allows.

\begin{figure}[ht]
\centering\begin{tikzpicture}[x=1cm,y=1cm,>=Latex,
 every node/.style={font=\small}]
 \draw[->,gray!65] (-4,0)--(4,0) node[right] {$\Re\tau_v$};
 \draw[->,gray!65] (0,-2.3)--(0,.65) node[above] {$\Im\tau_v$};
 \draw[very thick,blue!65!black,->] (-3,0)--(3,0);
 \draw[very thick,orange!80!black,->] (-3,-1.6)--(3,-1.6);
 \draw[dashed,thick] (-3,0)--(-3,-1.6);
 \draw[dashed,thick] (3,-1.6)--(3,0);
 \node[above] at(-3,0){$-\rho$};
 \node[above] at(3,0){$\rho$};
 \node[above left] at(0,0){$0$};
 \fill (1.0,-1.6) circle(2pt);
 \node[below right] at(1,-1.6){$\tau_{\boldsymbol\ell,v}$};
 \node[left] at(-3.1,-.8){connector};
 \node[right] at(3.1,-.8){connector};
 \node[below left] at(0,-1.6){$-v_v$};
\end{tikzpicture}

\caption{One coordinate of the finite angular deformation.
The real box has endpoints \(\pm\rho\), with \(\rho=1/8\);
its translated horizontal side passes through the critical point.
The two vertical sides are retained as connectors, with their
boundary orientations when comparing the horizontal integrals.
Here \(v_v=-\Im\tau_{\boldsymbol\ell,v}>0\).
Other coordinates can have intermediate downward displacements,
so their negative Fourier terms are kept in
\eqref{cta:connector}. The earlier infinite-charge estimate
uses real or translated-grid spectators, before this deformation.}
\label{cta:contour-figure}
\end{figure}

\subsection{Matching the geometric branch}
The analytic action is exactly the action of
\cite[Theorem~1.1]{Geometry}, not just a function with its
complete Hessian. Its derivative \eqref{cta:psi-definition}
agrees on every edge with that theorem, and the constants agree:
\(2G_0=v_{\rm oct}/\pi\), where \(v_{\rm oct}\) is the regular
ideal octahedron volume. This last equality also follows directly
from the Lobachevsky integral at \(\pi/4\).

Write \(n_v=\epsilon_v-2p_v\), with the retained parity
\(\epsilon_v\in\{-1,0,1\}\) in each actual disk frame.
The geometric theorem fills the original primitive slope
\(\mu_v+p_v\lambda_v\), with periods
\[
 U_v=-2\pi\Psi_{\tau_v}-\ii\epsilon_v\tau_v,\qquad
 V_v=2\ii\tau_v.
\]
Its unit-charge equation is
\(\Psi_{\tau_v}+\ii n_v\tau_v/(2\pi)=\ii\), precisely
\eqref{cta:fixedpoint} for \(\boldsymbol\ell=\mathbf1\).
The summed condition implies
\(\sqrt{d_v^2+n_v^2}\ge2^{25/2}>630\).
Furthermore \eqref{cta:tau-formula} gives
\[
 |\tau_{*,v}|\le
 \frac{2\pi}{(1-\delta)\sqrt{d_v^2+n_v^2}}<10^{-2}.
\]
Thus our critical point lies in the geometric theorem's
uniqueness chart. Its original filled metric is smooth,
complete and finite-volume, including zero powers at high-degree
vertices; every main cusp stays complete. Its absolute
fundamental-chain calculation gives
\(\Vol(S^3\setminus K_{\mathbf n})=2\pi\Rea\calC(\mathbf1)\).
There is no appeal here to a long-twist theorem for a zero
power, or to a formal exponentiated slope without metric completion.
Taking logarithms of the already nonzero
\eqref{cta:leading} proves Theorem~\ref{cta:main}.

\subsection{The centered main-meridian insertion}
The finite scalar ledger of \cite[Proposition~8.1]{Meridian}
separates exactly \(A_N(s)\) from the index-dependent factors
\(H_N/W_k\), the shifted crossing grids and the single closing
factor \(H_N(s)\). In the first-charge coefficient every grid
parameter has disappeared because \(c_{1,N}(\sigma)=1\).
The chirp has no remaining \(\sigma\)-dependence once its
recorded scalar is extracted.

It remains to integrate the second line of
\eqref{cta:packet-expansion} on the same joint contour.
For any fixed holomorphic insertion \(h(\boldsymbol\tau)\),
the normalized angular Gaussian evaluates it as
\(h(\boldsymbol\tau_*)+O(N^{-1})\): its linear Gaussian
term vanishes, and the quadratic term and its interaction
with the cubic phase contribute \(O(N^{-1})\).
The expansion to the next order and the fixed losses above
make this statement uniform on \(|s|\le R\).
Apply it to
\(-\pi s^2\sum_e\cot(\pi x_e(\boldsymbol\tau))\),
which already has the factor \(N^{-1}\).
The common first saddle correction cancels upon division
by \(F_N(1)\). The result is \eqref{cta:profile-formula}
with a uniform \(O_R(N^{-2})\) remainder.
In particular the wire values are those at the filled
critical point, not the complete values \(1/4\).

Both \(A_N\) and \(H_N\) are nonzero and bounded away from
zero on the fixed disk for large \(N\).
The ratio in \eqref{cta:profile-formula} tends uniformly
to one. This proves uniform eventual nonvanishing and
the remaining assertions of Theorem~\ref{cta:profile}.
No division by an unproved nonzero finite invariant was
used to establish the leading term.

\section{Scope, dependencies and reproducibility}

The condition is a theorem about every fixed literal marked
tree satisfying \eqref{cta:sum}. It is not uniform in a tree
growing with \(N\), and it is not a condition solely on the
smallest crossing power. For example, in a star with \(d\)
terminal vertices, central power zero and terminal powers \(n\),
\[
 \mathfrak S=\frac1d+\frac{d}{1+n^2}.
\]
Taking \(d\ge2^{27}\) and \(|n|\ge\sqrt{2^{27}d}\) gives
\(\mathfrak S\le2^{-26}\). Knot applicability still requires
the literal endpoint condition of Section~\ref{cta:intro};
this arithmetic example does not supply that condition by itself.
The proof uses neither the separate directed-message criterion
nor the more general strict Conway grammar with bare junctions.

The mathematical dependencies have different jobs.
The finite meridian calculus supplies the exact original
invariant and all repeated-root data.
The factorial lemma supplies one normalized holomorphic
interpolant, and its application to this finite edge sum is
proved in Section~\ref{cta:edge}.
Sections~\ref{cta:angular}--\ref{cta:Fourier} prove the joint
summed-condition analytic estimate, including the entire
original source and Fourier period.
Only after noncancellation is proved is the classical
circle-tree theorem used to identify the action with metric
volume. No quantum asymptotic is an input to that geometric
theorem.

The residue and Gaussian steps belong to the classical
contour method; \cite[Sections~2.4(iii)--(iv)]{DLMF} gives
background. Their contours and all losses relevant to this
finite invariant have been specified above.
From the repository root the paper-local reproduction command is
\begin{verbatim}
python3 papers/circle-tree-asymptotics/check.py \
    --output /tmp/circle-tree-check.json
\end{verbatim}
It records source hashes and reuses the existing edge script's four
exact symbolic identities and finite complex-product samples.
It also checks the numerical budgets by exact rational interval
arithmetic, using convergent series bounds for \(\pi\), \(\sinh\)
and \(\cosh\). Negative controls reject altered budgets.
Those finite calculations do not prove a uniform strip estimate,
the charge tails or a volume theorem; the corresponding
proofs are given here. The script and its recorded result are
linked in the accompanying repository.

\begin{thebibliography}{9}
\bibitem{Meridian}
\emph{Exact Newton calculus for tangles with a moving meridian},
companion paper, Proposition~6.1, Theorem~7.1,
and Proposition~8.1.
\href{https://sashakolpakov.github.io/murakamis_dream/papers/meridian-tangles.pdf}
{Companion PDF}.

\bibitem{Geometry}
\emph{Positive octahedra and primitive fillings for trees of crossing circles},
companion paper, Theorem~1.1.
\href{https://sashakolpakov.github.io/murakamis_dream/papers/circle-tree-geometry.pdf}
{Companion PDF}.

\bibitem{ConwayAnalysis}
\emph{Holomorphic source asymptotics and volume limits for Conway closures},
companion paper, Lemma~2.1.
\href{https://sashakolpakov.github.io/murakamis_dream/papers/conway-asymptotics.pdf}
{Companion PDF}.

\bibitem{Oscillatory}
\emph{Oscillatory filling sums with independently long twists},
companion paper.
\href{https://sashakolpakov.github.io/murakamis_dream/papers/oscillatory-fillings.pdf}
{Companion PDF}.

\bibitem{DLMF}
NIST Digital Library of Mathematical Functions,
\href{https://dlmf.nist.gov/5.11}{Section~5.11, Gamma-function asymptotics},
and \href{https://dlmf.nist.gov/2.4}{Section~2.4, Contour integrals}.
\end{thebibliography}
